\documentclass[12pt]{article}
\usepackage[T1]{fontenc}
\usepackage[utf8]{inputenc}
\usepackage{lmodern}
\usepackage{textcomp}
\usepackage{enumitem}
\usepackage{geometry}
\geometry{margin=1in}
\begin{document}
\begin{center}
Answer Key - Political Science - Undergraduate Level Assessment 18 \
\small (Answers are ordered exactly as in the question paper.)
\end{center}

\section*{Multiple Choice}
\begin{enumerate}[label=\arabic*.]
\item \textbf{Answer:} C
\\ \textit{Options:} A) bread and butter.; B) war and peace.; C) guns and butter.; D) bombs and books.
\\ \textit{Note:} The phrase "guns and butter" encapsulates the economic and political trade-off between allocating resources to military defense (guns) versus social programs and welfare (butter), highlighting the opportunity costs involved in government spending decisions.
\item \textbf{Answer:} B
\\ \textit{Options:} A) peer pressure; B) economic issues; C) religious beliefs; D) pressure from their employers
\\ \textit{Note:} Individuals often align with political parties based on their personal experiences and beliefs regarding economic issues, which can differ significantly from their parents' perspectives, leading them to seek affiliation with parties that better represent their views on economic policy.
\item \textbf{Answer:} B
\\ \textit{Options:} A) A line-item veto; B) A pocket veto; C) An adjournment; D) A writ of certiorari
\\ \textit{Note:} A pocket veto occurs when the president takes no action on a bill for ten days while Congress is in recess, effectively preventing the bill from becoming law without issuing a formal veto.
\item \textbf{Answer:} D
\\ \textit{Options:} A) A growing effort by governments to prevent the spread of weapons, especially of WMDs.; B) A growth in volume of the arms trade.; C) A change in the nature of the defence trade, linked to the changing nature of conflict.; D) All of these options.
\\ \textit{Note:} The correct answer is "All of these options" because the global defense trade has indeed exhibited multiple interrelated trends since the end of the Cold War, including increased privatization, the rise of emerging powers, and the proliferation of advanced military technology.
\item \textbf{Answer:} A
\\ \textit{Options:} A) Constitutional amendments; B) Presidential executive orders; C) Laws passed by Congress; D) Laws passed by state legislatures
\\ \textit{Note:} The Supreme Court does not have the power to override constitutional amendments because such amendments are a fundamental aspect of the Constitution that require a rigorous legislative process for alteration, ensuring that they are beyond the reach of judicial authority.
\end{enumerate}

\section*{True / False}
\begin{enumerate}[label=\arabic*.]
\setcounter{enumi}{5}
\item \textbf{Answer:} False
\\ \textit{Options:} A) True; B) False
\\ \textit{Note:} The statement is false because, in the event that no presidential candidate receives a majority of the electoral college votes, the election is decided by the House of Representatives, which selects the president from the top three candidates, while the Senate chooses the vice president.
\item \textbf{Answer:} True
\\ \textit{Options:} A) True; B) False
\\ \textit{Note:} Democratic enlargement is accurately defined as the strategy aimed at promoting democracy and free market economies worldwide, reflecting the belief that these systems contribute to global stability and peace.
\item \textbf{Answer:} True
\\ \textit{Options:} A) True; B) False
\\ \textit{Note:} The Constitution does not include any specific language regarding judicial review; instead, this power was established by the Supreme Court in the landmark case Marbury v. Madison in 1803.
\item \textbf{Answer:} False
\\ \textit{Options:} A) True; B) False
\\ \textit{Note:} Motor voter laws allow individuals to register to vote while applying for or renewing a driver's license, but they do not mandate that all voters must possess a driver's license to participate in elections, as other forms of identification or registration options are also accepted.
\item \textbf{Answer:} True
\\ \textit{Options:} A) True; B) False
\\ \textit{Note:} Federal budget entitlements are considered mandatory spending programs that automatically provide benefits to eligible individuals or groups, such as Social Security and Medicare, without requiring annual approval through an appropriations bill.
\end{enumerate}

\section*{Long Form Response}
\begin{enumerate}[label=\arabic*.]
\setcounter{enumi}{10}
\item \textbf{Answer:} The long-term implications of environmental security on global policy are profound, as the persistence of environmental issues suggests that the importance of incorporating environmental security into policy frameworks will only grow. Multilateral action will become essential for addressing common security concerns, fostering social organization that prioritizes ecological sustainability. However, current national security practices often hinder the pursuit of environmental security, as they tend to focus on militarization and state-centric approaches rather than collaborative solutions. This creates a paradox where national security priorities can exacerbate environmental degradation, ultimately threatening both human security and global stability. Therefore, reconciling national security with environmental imperatives will be crucial for developing effective and sustainable policy responses in the future.
\\ \textit{Note:} correct because it highlights the increasing necessity of integrating environmental security into global policy due to persistent environmental challenges and emphasizes that current national security practices can impede this integration.
\item \textbf{Answer:} The financial support provided by U.S. banks to the Allies significantly influenced America's declared neutrality in World War I between 1915 and 1917. American banks, particularly J.P. Morgan \& Co., extended substantial loans to the Allies, amounting to billions of dollars, while providing minimal financial assistance to Germany and Austria- Hungary. This overwhelming economic support created a vested interest for the United States in the success of the Allied powers, as a victory for the Allies would ensure repayment of loans and stabilize U.S. economic interests. Consequently, although the U.S. maintained a neutral stance publicly, its financial entanglements with the Allies fostered a pro-Allied sentiment among policymakers and the public, ultimately eroding the commitment to neutrality as the war progressed. This financial dynamic demonstrated how economic interests can shape national policy and influence the trajectory of foreign relations.
\\ \textit{Note:} correct because the substantial financial loans from U.S. banks to the Allies created a strong economic dependency and vested interest for America, ultimately undermining its declared neutrality as the nation's prosperity became linked to the success of the Allied powers in the war.
\end{enumerate}

\vfill
\noindent\textit{End of Answer Key}
\end{document}